\documentclass[11pt]{article}

\usepackage[margin=1in]{geometry}
\usepackage{amsmath,amssymb,bm}
\usepackage{enumitem}
\usepackage{hyperref}

\newcommand{\dd}{\mathrm{d}}
\newcommand{\Dt}{\Delta t}
\newcommand{\lapS}{\nabla_{\mathbb{S}}^{2}}
\newcommand{\jacS}{J_{\mathbb{S}}}
\newcommand{\Ylm}{Y_{\ell m}}
\newcommand{\chiM}{\chi}
\newcommand{\sponge}{\sigma_{\mathrm{sp}}}

\title{Two-Layer QG Feasibility Model on the Sphere via Dinosaur}
\author{NHQG / Dinosaur spike}
\date{2026-05-22}

\begin{document}
\maketitle

\section{Purpose}

This document defines the equations and numerical method for the first
Dinosaur/NeuralGCM feasibility spike.  The goal is not yet the full
nonhydrostatic QG equation.  The goal is to test whether Dinosaur's
global spherical-harmonic transforms, differential operators, and sharding
machinery can support a custom QG-like model with broad high-latitude and
midlatitude dynamics.

The first dynamical test is a two-layer Phillips-style QG model.  This is
more informative than barotropic vorticity because it includes baroclinic
coupling and deformation-radius physics, the part of QG that becomes most
delicate when a formally high-latitude/asymptotic model is embedded on the
full sphere.

\section{Geometry and representation}

We work on the unit sphere unless stated otherwise.  Latitude is
$\varphi \in [-\pi/2,\pi/2]$, longitude is $\lambda$, and the spherical
area element is
\begin{equation}
  \dd A = \cos\varphi\,\dd\varphi\,\dd\lambda.
\end{equation}
The Coriolis parameter is
\begin{equation}
  f(\varphi) = 2\Omega\sin\varphi.
\end{equation}

Fields are represented globally in spherical harmonics:
\begin{equation}
  a(\lambda,\varphi,t)
  = \sum_{\ell=0}^{L_{\max}}\sum_{m=-\ell}^{\ell}
  \widehat{a}_{\ell m}(t)\,\Ylm(\lambda,\varphi),
\end{equation}
with the exact coefficient layout and normalization inherited from
Dinosaur.  The spike must record that convention in code reviews once the
API is inspected.

The spherical Laplacian satisfies
\begin{equation}
  \lapS \Ylm = -\ell(\ell+1)\Ylm
\end{equation}
on the unit sphere.  If Dinosaur uses dimensional radius $a$, replace
$\ell(\ell+1)$ by $\ell(\ell+1)/a^{2}$.

\section{Two-layer QG variables}

The prognostic variables are layer potential-vorticity anomalies
$q_1'(\lambda,\varphi,t)$ and $q_2'(\lambda,\varphi,t)$.
Layer streamfunctions are $\psi_1$ and $\psi_2$.

The inversion from streamfunction to PV anomaly is
\begin{align}
  q_1' &= \lapS \psi_1 + F_1(\psi_2-\psi_1),\\
  q_2' &= \lapS \psi_2 + F_2(\psi_1-\psi_2).
  \label{eq:qg-inversion}
\end{align}
For the first implementation, $F_1,F_2$ are constants chosen to represent
the target high-latitude/midlatitude deformation scale.  This makes the
inversion block diagonal in spherical harmonic degree $\ell$.

For equatorial-regularization tests we also allow latitude-dependent effective
couplings
\begin{align}
  F_i(\varphi) &= F_i^\ast
  \frac{\sin^2\varphi + s_{\min}^2}
       {\sin^2\varphi_\ast + s_{\min}^2},
  \label{eq:f-squared-floor}
\end{align}
where $\varphi_\ast$ is a reference latitude and $s_{\min}>0$ floors the
Coriolis factor.  This is the ``$f^2$ floor'' profile.  A reciprocal diagnostic
profile replaces the final ratio by its inverse.  Once $F_i$ varies with
latitude, Eq.~\eqref{eq:qg-inversion} is no longer diagonal in spherical
harmonic degree, so the spike uses dense fixed-$m$ linear inversions before
attempting nonlinear time stepping.

\section{Spectral inversion}

Let $L_\ell = \ell(\ell+1)$ on the unit sphere.  Since
$\lapS \to -L_\ell$, Eq.~\eqref{eq:qg-inversion} gives, for each
$(\ell,m)$,
\begin{equation}
  \begin{bmatrix}
    \widehat{q_1'}\\
    \widehat{q_2'}
  \end{bmatrix}_{\ell m}
  =
  \begin{bmatrix}
    -L_\ell - F_1 & F_1\\
    F_2 & -L_\ell - F_2
  \end{bmatrix}
  \begin{bmatrix}
    \widehat{\psi_1}\\
    \widehat{\psi_2}
  \end{bmatrix}_{\ell m}.
  \label{eq:spectral-inversion}
\end{equation}
The $(\ell,m)=(0,0)$ barotropic null mode is gauge-fixed by setting the
mean streamfunction to zero.  Operationally, the inversion sets
$\widehat{\psi_1}_{00}=\widehat{\psi_2}_{00}=0$ and removes any mean PV
component that is not dynamically meaningful.

\section{Background zonal shear}

To test Phillips-style baroclinic instability, the implemented model may
evolve perturbations about a prescribed zonal base state
$\Psi_i^0(\varphi)$.  The first base state is the smooth global profile
\begin{equation}
  \Psi_i^0(\varphi) = -U_i\sin\varphi,
\end{equation}
which gives zonal velocity $u_i^0 = U_i\cos\varphi$ on the unit sphere.
Exact spherical linear checks show that this solid-body profile is a neutral
control rather than a useful Phillips-instability case.  A second regular
curved-shear profile is therefore
\begin{equation}
  \Psi_i^0(\varphi)
  =
  -U_i\left(\sin\varphi + a_3\sin^3\varphi\right),
  \qquad a_3 = 0.75 \text{ by default}.
\end{equation}
The curved profile remains polar-regular because its zonal velocity is
proportional to $\cos\varphi$, but it has nontrivial meridional PV curvature
and supports unstable two-layer modes in the linear spectrum.
We parameterize
\begin{align}
  U_1 &= U_{\mathrm{bt}} + \frac{1}{2}U_{\mathrm{bc}},\\
  U_2 &= U_{\mathrm{bt}} - \frac{1}{2}U_{\mathrm{bc}},
\end{align}
where $U_{\mathrm{bc}}$ is the layer shear.

The base PVs are
\begin{align}
  Q_1^0 &= \lapS \Psi_1^0 + F_1(\Psi_2^0-\Psi_1^0) + f,\\
  Q_2^0 &= \lapS \Psi_2^0 + F_2(\Psi_1^0-\Psi_2^0) + f.
\end{align}
The prognostic state remains the perturbation PV anomaly $q_i'$.

\section{Spherical Jacobian}

The advective Jacobian on the sphere is
\begin{equation}
  \jacS(a,b)
  =
  \frac{1}{\cos\varphi}
  \left(
    \frac{\partial a}{\partial \lambda}
    \frac{\partial b}{\partial \varphi}
    -
    \frac{\partial a}{\partial \varphi}
    \frac{\partial b}{\partial \lambda}
  \right),
\end{equation}
up to the sign convention used by Dinosaur's vorticity/velocity helpers.
The spike must verify this sign convention using a manufactured field test.

\section{Masked/tapered evolution}

The broad target is not a hard polar cap.  We introduce a smooth
Tukey-style latitude envelope $\chiM(\varphi)$:
\begin{equation}
  0 \le \chiM(\varphi) \le 1.
\end{equation}
For a southern-hemisphere configuration, the taper may begin equatorward
of approximately $30^\circ$S, with a broad transition toward the invalid
or less-trusted low-latitude reservoir.  The exact plateau and taper
locations are experimental parameters.

The perturbation tendency model is
\begin{align}
  \partial_t q_1'
  &=
  -\chiM\left[
    \jacS(\psi_1',q_1'+Q_1^0)
    + \jacS(\Psi_1^0,q_1')
  \right]
  - \sponge(1-\chiM)\,q_1'
  - \mathcal{D}q_1'
  + \chiM\,\mathcal{F}_1,\\
  \partial_t q_2'
  &=
  -\chiM\left[
    \jacS(\psi_2',q_2'+Q_2^0)
    + \jacS(\Psi_2^0,q_2')
  \right]
  - \sponge(1-\chiM)\,q_2'
  - \mathcal{D}q_2'
  + \chiM\,\mathcal{F}_2.
  \label{eq:tendency}
\end{align}
Here $\mathcal{D}$ is hyperdiffusion or spectral filtering and
$\mathcal{F}_{1,2}$ are optional relaxation/forcing terms.

This is a deliberate first model, not a conservation theorem.  Multiplying
nonlinear tendencies by $\chiM$ will generally break exact PV/enstrophy
conservation.  The spike must diagnose whether the resulting spectral
ringing/leakage is acceptable.  Alternative variants include applying
$\chiM$ only to forcing and sponge terms, leaving the advective tendency
unmasked.

\section{Time stepping}

The first implementation should use a simple explicit RK4 or SSPRK scheme:
\begin{equation}
  q'^{n+1} = \mathrm{RK}(q'^n,\Dt).
\end{equation}
This keeps the feasibility spike focused on transforms, inversion, masking,
and sharding.  IMEX machinery is not needed until the model graduates toward
NHQGE vertical coupling.

\section{Numerical method checklist}

\begin{enumerate}[leftmargin=2em]
  \item Build or import Dinosaur coordinate and spherical-harmonic transform
  objects.
  \item Verify spectral-to-nodal-to-spectral roundtrip for scalar fields.
  \item Verify $\lapS \Ylm = -\ell(\ell+1)\Ylm$ in Dinosaur's convention.
  \item Implement Eq.~\eqref{eq:spectral-inversion} and verify recovery of
  manufactured $\psi_1,\psi_2$.
  \item Implement $\jacS(\psi_i,Q_i)$ using Dinosaur derivative/velocity
  helpers where possible.
  \item Add $\chiM(\varphi)$, sponge, forcing, and diagnostics.
  \item Benchmark the same step function on one GPU and with Dinosaur-style
  sharding over multiple GPUs.
\end{enumerate}

\section{Diagnostics}

The spike should report at least:
\begin{itemize}[leftmargin=2em]
  \item spectral roundtrip residuals,
  \item inversion residuals for Eq.~\eqref{eq:spectral-inversion},
  \item maximum and spectrum of masked tendency fields,
  \item energy/enstrophy-like diagnostics inside the trusted window,
  \item leakage diagnostics in the sponge/low-latitude reservoir,
  \item wall time per step and transform count.
\end{itemize}

\end{document}
